% Generated draft: requires mathematical and source review.
\documentclass[11pt]{article}
\usepackage[margin=25mm]{geometry}
\usepackage{fontspec}
\setmainfont{Latin Modern Roman}
\usepackage{amsmath,amssymb,mathtools,unicode-math}
\setmathfont{Latin Modern Math}
\usepackage{enumitem,needspace,xurl,hyperref}
\setlength{\parindent}{0pt}
\setlength{\parskip}{5pt}
\setlength{\emergencystretch}{3em}
\newcommand{\N}{\mathbb{N}}
\newcommand{\Z}{\mathbb{Z}}
\newcommand{\Q}{\mathbb{Q}}
\newcommand{\R}{\mathbb{R}}
\newcommand{\C}{\mathbb{C}}
\newcommand{\F}{\mathbb{F}}
\newcommand{\A}{\mathbb{A}}
\newcommand{\PP}{\mathbb{P}}
\newcommand{\E}{\mathbb{E}}
\newcommand{\eps}{\varepsilon}
\newcommand{\dd}{\,\mathrm{d}}
\begin{document}
\section*{UnsolvedMath ID: 11000141}
\textbf{Source problem code:} catalogue entry
\subsection{Definitions and mathematical statement}
Let \(T_N(F)\) denote the group called the level-\(N\) Torelli group of the surface \(F\) in the source. The Level-\(N\) Torelli Franchetta Problem asks to compute its second group cohomology,
\[
 H^2(T_N(F);A),
\]
with the coefficient module \(A\) and action fixed by the missing conventions. Integral trivial coefficients are a usual interpretation, but are not printed here. This is cohomology, not second homology or merely a rational Betti number.

The definition of level-\(N\) Torelli group is material and absent. Ordinary Torelli acts trivially on integral first homology, so intersecting it with a principal homological level subgroup would leave it unchanged; that cannot be assumed to specify an intended varying family. The source's actual level construction, admissible \(N\), genus and marking conventions must be supplied.

The accompanying text describes an ideal-cell decomposition of \(\overline{M}(F)\) compatible with Torelli-space decompositions: the dual fatgraph, a graph with cyclic orders of incident half-edges, has both a homology marking in the sense of [48] and screen data. It says this yields Deligne--Mumford-type boundaries of each Torelli space with ideal-cell decompositions, and proposes studying the topology of these bordifications to approach further problems. The definitions of screens, homology markings and these particular bordifications, as well as the following class of problems, are not supplied. No specific compactification or Chow-theoretic reformulation is substituted for the group-cohomology question.
\subsection{Short English statement}
Compute second cohomology of the source's level-\(N\) Torelli group, using the intended level and coefficient definitions; the proposed fatgraph/screen bordifications are accompanying motivation.
\subsection{Sources}
\begin{itemize}
\item \emph{Problems on Mapping Class Groups and Related Topics}, edited by Farb, Problem 11, ``Level\(N\) Torelli Franchetta Problem,'' supplied PDF page 121. \url{https://www.math.uchicago.edu/~farb/papers/mcgbook.pdf}. Recorded access 2026-07-29; extraction pdf-font-structure.
\item Source [48] for the homology-marking convention; its bibliography entry and the preceding ideal-cell/screen constructions are absent.
\item Local UnsolvedMath registry, \texttt{data/unsolvedmath/problems.json}, permanent ID \(11000141\), exact source code \texttt{catalogue entry}; supplied OPEN-TRIAGE summary mentions related universal-family Chow/cohomology work without exact citations.
\end{itemize}
\end{document}
